% ============================================================
% Chapter 4 -- Methodological Approach and Design Choices
% File: chapters/chapter4.tex
%
% Preamble reminder for the main .tex (Rapport_Stage.tex):
%   Nothing to add. tnreport.cls already loads booktabs and tikz with
%   the positioning/arrows.meta/fit libraries used by the Section 4.2
%   figure, and listings for code excerpts. Long \texttt{} identifiers
%   are wrapped with manual \allowbreak instead of a seqsplit package,
%   per the project convention (see claude/memoire-reference.md, S2).
%
% Included from the main file as:
%   \input{chapters/chapter4}
%   \cleardoublepage{}
% ============================================================

\chapter{Methodological Approach and Design Choices}
\label{chap:method}
\emergencystretch=3em\relax

This chapter states the approach that shaped every technical choice
described later in this report. It sits between the problem statement
developed in Chapter~\ref{chap:context} and the two chapters that describe
how the platform was actually built, Chapter~\ref{chap:scoring-engine} on
the scoring engine and Chapter~\ref{chap:platform} on the surrounding
architecture. It appears here in the reading order, but I wrote it after
those two chapters were already drafted, once the implementation had
settled enough to state the reasoning behind it clearly, rather than as it
was being discovered along the way.

The chapter is organized in four parts. Section~\ref{sec:working-method}
describes how the project was run sprint by sprint. Section~\ref{sec:core-idea}
states the central design idea of the platform, the separation between
deterministic scoring and AI augmentation, and why it answers the problem
raised in the state of the art. Section~\ref{sec:principles} lists the five
principles that were fixed early and never revisited. Section~\ref{sec:design-choices}
details the main arbitrations made against that background, together with
the alternatives that were considered and rejected.

\section{Working method}
\label{sec:working-method}

The project was organized as a sequence of bounded sprints, each defined by
a written brief and closed by an explicit Definition of Done. Sprint~0, for
instance, listed eight verifiable conditions to meet before the sprint
could be considered closed, from every scoring engine test passing to a
clean \texttt{git status}, and stated explicitly what was left out of scope
for that sprint. Closing a sprint meant checking that list item by item,
instead of declaring the work finished by feel.

This discipline was paired with a documentation-first habit, built around
two documents. \texttt{CLAUDE.md} is the framing document for
the whole project: it states the five inviolable principles discussed in
Section~\ref{sec:principles}, the working method itself, down to the rule
that an ambiguous sprint brief should be clarified before coding rather
than settled by guessing, and a numbered list of anti-patterns to avoid,
reproduced in Table~\ref{tab:anti-patterns} so that later chapters can
refer to them by number instead of restating them.
\texttt{ARCHITECTURE.\allowbreak md} is an append-only, dated log of the
architectural decisions taken sprint after sprint, each entry recording
both the decision and its justification. Both documents were written as
the decisions were being made, which is why most of the justifications
given in Section~\ref{sec:design-choices} are drawn directly from that log
rather than reconstructed from memory for this report.

\begin{table}[htbp]
  \centering
  \begin{tabular}{@{}cp{0.82\textwidth}@{}}
    \toprule
    \# & Anti-pattern \\
    \midrule
    1 & Reimplementing the scoring engine, its loader, or the backend layers instead of extending them. \\
    2 & Introducing LangChain, LlamaIndex, or any AI framework on top of Ollama. \\
    3 & Letting the AI compute or modify a score. \\
    4 & Hardcoding referential text in Python or TypeScript instead of loading it. \\
    5 & Reintroducing anonymous access or bypassing the centralized auth layer. \\
    6 & Splitting the stack into microservices. \\
    7 & Relying on implicit SQLAlchemy lazy loading where it can hide N+1 queries. \\
    8 & Committing \texttt{.env} files or other secrets. \\
    9 & Writing a frontend component longer than 200 lines. \\
    10 & Reordering the scoring pipeline, or changing the referential's structure, without bumping its version. \\
    \bottomrule
  \end{tabular}
  \caption{Anti-patterns to avoid, as fixed in \texttt{CLAUDE.md}}
  \label{tab:anti-patterns}
\end{table}

Development followed a test-driven approach wherever the scoring logic was
involved. The engine carries seventy tests, including golden scenarios
that double as an executable specification of the scoring rules, and it
went through an audit that fed corrections back into the code before the
engine was integrated into the backend. The audit was carried out by a
separate, freshly initialized review with no prior exposure to the
implementation, which is what makes it independent, rather than by a
third-party reviewer. Its testing strategy and the audit itself are
detailed in Chapter~\ref{chap:scoring-engine}, which this chapter does not
repeat.

Tooling choices, such as the backend and frontend stacks, follow from this
method more than they define it, and are presented where they matter, in
Chapter~\ref{chap:platform}.

\section{The core idea: separating deterministic scoring from AI augmentation}
\label{sec:core-idea}

Two problems motivate the central design decision described in this
section. The first is auditability: because the deliverable is meant to
give a consultant a defensible, reproducible score to bring in front of a
client, the computation itself cannot depend on anything that is not fully
reproducible. The second is hallucination risk: Chapter~\ref{chap:state-of-the-art}
identifies the tendency of large language models to produce plausible but
false content as the main obstacle to their adoption in sensitive domains
such as cyber crisis management. Both problems point to the same
conclusion: whatever a language model contributes to this platform, it
must never be allowed to touch the number that ends up in a client's
report.

The response is to split the system into two stages that never share write
access to the score. A pure, deterministic engine consumes the assessment
inputs together with the framework defined in the referential, and
produces a fully computed result with no AI involvement at any point,
following the determinism principle detailed in Section~\ref{principle:determinism}.
Only once that result exists does a second, local language model read it
to produce narrative text, an executive synthesis and a set of
recommendations, and nothing else: the referential itself states this
explicitly, marking the AI's scoring role as none. Splitting the pipeline
this way addresses both problems at once. Auditability holds because the
score is always reproducible from the raw inputs and the referential
version alone, without ever invoking a non-deterministic component. The
hallucination risk is contained rather than eliminated: the language model
can still write something inaccurate, but the blast radius of that mistake
is limited to a downstream text field that can be checked, regenerated, or
replaced by a static fallback, and it can never reach back into the score
itself. How that containment is implemented, the prompt guardrails and the
fallback mechanism, belongs to Chapter~\ref{chap:platform} and this
chapter does not repeat it.

Figure~\ref{fig:scoring-ai-separation} illustrates the resulting pipeline.
Assessment inputs flow into the deterministic engine, which produces the
\texttt{Assessment\allowbreak Result}; that result is then read, never written, by the local
language model, which produces narrative output only. The diagram
deliberately does not draw an arrow running back from the language model
to the engine or to the score: nothing downstream of the score is allowed
to influence it.

\begin{figure}[htbp]
  \centering
  \begin{tikzpicture}[
      node distance=0.9cm and 1.0cm,
      box/.style={draw, rounded corners, align=center, minimum width=2.9cm, minimum height=1.05cm, font=\small},
      arr/.style={-{Straight Barb[angle=60:2mm]}, thick},
      lbl/.style={font=\scriptsize, align=center, fill=white, inner sep=1pt},
      stagelabel/.style={font=\footnotesize\itshape}
    ]
    \node[box] (inputs) {Assessment\\ inputs};
    \node[box, right=of inputs] (engine) {Deterministic\\ scoring engine\\ \tiny (pure Python, P1)};
    \node[box, right=of engine] (result) {AssessmentResult\\ \tiny (not stored)};

    \node[box, below=1.7cm of engine] (llm) {Local LLM\\ \tiny (Ollama)};
    \node[box, right=of llm] (narrative) {Narrative\\ output\\ \tiny (text only)};

    \draw[arr] (inputs) -- (engine);
    \draw[arr] (engine) -- (result);
    \draw[arr] (result.south) .. controls +(0,-0.8) and +(0,0.8) .. node[lbl, right] {read-only} (llm.north);
    \draw[arr] (llm) -- (narrative);

    \node[draw, dashed, rounded corners, fit=(inputs)(engine)(result), inner sep=8pt, label={[stagelabel]above:Deterministic stage}] {};
    \node[draw, dashed, rounded corners, fit=(llm)(narrative), inner sep=8pt, label={[stagelabel]below:AI augmentation stage}] {};
  \end{tikzpicture}
  \caption{Separation between the deterministic scoring engine and the AI augmentation layer}
  \label{fig:scoring-ai-separation}
\end{figure}

\section{Five inviolable principles}
\label{sec:principles}

The five principles listed in this section were fixed at the start of the
project, stated together in a single reference document, and never
revisited afterward. Table~\ref{tab:principles-summary} summarizes them;
each is then justified in turn, together with a note on where it is
enforced and where its consequences are detailed further in this report.

\begin{table}[htbp]
  \centering
  \begin{tabular}{@{}p{0.32\textwidth}p{0.58\textwidth}@{}}
    \toprule
    Principle & Rationale \\
    \midrule
    P1 -- Deterministic scoring & The engine alone computes the score; the AI never touches it, so the result is always reproducible and auditable. \\
    P2 -- Local-first & All language model inference runs on a local Ollama instance; no assessment data ever reaches an external API. \\
    P3 -- Secure by design & Data sensitivity drives the architecture, from encryption at rest to input validation and a single access-control module. \\
    P4 -- Single source of truth & \texttt{referential.\allowbreak yaml} is the only specification of the framework; no code ever hardcodes a question, a weight, or a rule. \\
    P5 -- Auth schema from day one & The data model carries users, organizations, and roles from the first sprint, so real authentication later replaces a stub without a refactor. \\
    \bottomrule
  \end{tabular}
  \caption{The five inviolable principles}
  \label{tab:principles-summary}
\end{table}

\subsection{P1: Determinism}
\label{principle:determinism}

Given the same answers, the same evidence levels, and the same referential
version, the scoring engine always returns the same score. The AI plays no
part in this computation at any stage; it only operates downstream on a
result that is already final. The guarantee matters beyond internal
consistency: it is what lets a consultant defend a client's number in
front of that same client, possibly months later, using the same rules
against a different set of answers. The engine enforces this by
construction, and a dedicated test compares two independent computations
on the same inputs to catch any accidental regression. Section~\ref{sec:core-idea}
already relied on this principle to justify separating scoring from AI
augmentation; the engine's internals and the audit that verified the
guarantee in practice are described in Chapter~\ref{chap:scoring-engine}.

\subsection{P2: Local-first}
\label{principle:local-first}

All language model inference in this platform runs through a locally
hosted Ollama instance. No assessment data is ever sent to an external
API, whether OpenAI, Anthropic, or a cloud-hosted Mistral endpoint. This
is a hard constraint rather than a preference for cost or latency reasons:
the NIS2 directive discussed in Chapter~\ref{chap:context} pushes
organizations toward demonstrable sovereignty over the data they process,
and an assessment produced by this platform is, in effect, a detailed map
of an organization's cyber weaknesses. Sending that map to a third-party
API, however reputable, would undermine the very trust the platform is
meant to build. Because inference stays local, no separate anonymization
gateway is required to satisfy sovereignty; the remaining data-at-rest
handling is covered under Section~\ref{principle:secure-by-design}. The
consequence of this principle on model selection, and the alternative it
ruled out, are developed in Section~\ref{sec:design-choices}.

\subsection{P3: Secure by design}
\label{principle:secure-by-design}

The sensitivity of the data processed by the platform, effectively a
structured account of an organization's crisis-management gaps, drives
architectural choices beyond the local-first constraint of
Section~\ref{principle:local-first}. Concretely, this means the database
volume is encrypted at rest, free-text fields such as N/A justifications
are never written to logs in clear text, all traffic runs over HTTPS in
production, and every API boundary validates its input rather than
trusting the caller. The same principle extends to the AI layer: the
prompt sent to the local language model carries scores, best-practice
identifiers, the rationale behind any gate or penalty applied, and
referential text, but never an organization name, an identifier, or the
content of a justification field. Access control
follows the same logic, centralized in a single module rather than
re-implemented per route, closing the kind of broken access control that
OWASP lists as the most common web application weakness~\cite{owasp_top10_a01_2021}.
The concrete
mechanisms, encryption, logging discipline, and the access-control module,
are described where they are implemented, in Chapter~\ref{chap:platform}.

\subsection{P4: The referential as the single source of truth}
\label{principle:referential-truth}

The specification of the ENISA framework, its best practices, questions,
weights, maturity scale, evidence caps, foundational gates, and
consistency penalties, lives in a single YAML file at the root of the
repository. Neither the scoring engine nor the frontend ever hardcodes a
question, a weight, or a rule: the engine loads the file once through a
single loader, and the frontend fetches the same content from the API
rather than duplicating it. This matters for two reasons. First, it
removes an entire class of bugs where the code and the framework silently
drift apart, for instance a question rephrased in the interface but not in
the engine's validation. Second, it makes the framework itself an
auditable artifact: reviewing what the platform actually measures means
reading one file, not reconstructing it from scattered pieces of code. Any
change to the framework goes through that file and its version number,
never through a code change alone. How this loading and validation is
implemented is detailed in Chapter~\ref{chap:scoring-engine}, and how the
frontend consumes it without duplication in Chapter~\ref{chap:platform}.

\subsection{P5: Auth schema from day one}
\label{principle:auth-schema}

The data model has included users, organizations, roles, and the
corresponding foreign keys on assessments since the very first sprint,
well before any real authentication existed. Early on, the platform ran on
an anonymous session stub that satisfied the same interface a real
authenticated user would later have to satisfy. The principle paid off
exactly as intended: when real authentication was built, the stub was
replaced behind that same interface, and nothing downstream, including the
assessment ownership model and the role checks, had to be reworked. The
alternative, bolting the user and role concept onto the data model after
the fact, would have meant migrating existing assessments and every query
that touches them. The authentication flows that now sit behind this
schema are described in Chapter~\ref{chap:platform}.

\section{Key design choices and rejected alternatives}
\label{sec:design-choices}

Section~\ref{sec:principles} stated the five principles as constraints.
This section shows three of the arbitrations they produced in practice:
for each, I state the choice actually made, why the rejected alternative
was genuinely tempting, and which principle it would have compromised.
Table~\ref{tab:rejected-alternatives} summarizes the three below.

\subsection{Cloud LLM API rejected for data sovereignty}

A cloud-hosted large language model, whether from OpenAI, Anthropic, or
Mistral's own hosted offering, would have been the easier route for the AI
augmentation layer. These services expose considerably larger and more
capable models than anything that can run on commodity hardware, need no
local inference stack to install, monitor, or keep warm, and would likely
have produced better narrative quality than what actually runs locally.
That gap in capability is real, and Chapter~\ref{chap:platform} returns to
it directly. It was set aside anyway because of Section~\ref{principle:local-first}:
the NIS2 context described in Chapter~\ref{chap:context}, together with the
sensitivity of assessment data, effectively a map of an organization's
cyber weaknesses, rules out sending that data to a third party regardless
of the quality gained. The choice of local inference over a cloud API was
therefore imposed by the platform's subject matter rather than decided
purely on technical merit.

Once local inference was fixed, the sizing question reproduced a smaller
version of the same trade-off. The referential lists four larger
candidates for the local model, \texttt{mistral:7b-\allowbreak instruct},
\texttt{mixtral:8x7b}, \texttt{llama3.1:\allowbreak 8b}, and
\texttt{qwen2.5:\allowbreak 7b}, three of them in the 7 to 8 billion
parameter range and the fourth a mixture-of-experts model, none of which is
what the platform actually runs: the model deployed is
\texttt{qwen2.5:\allowbreak 3b}, smaller than any of them.
None of the larger candidates was actually benchmarked on the
project's hardware, only the deployed 3B model's generation time was
measured, at roughly a minute on the CPU-bound machine used for the demo.
The choice of the smaller model rested on the reasonable expectation that a
7B model would be slower still on the same hardware rather than on a
measured comparison, and \texttt{CLAUDE.md} itself defers that formal
benchmark to a later sprint. The resulting quality loss was judged
acceptable because the AI's task stays narrow, summarizing and enriching a
fixed set of referential recommendations for a short list of weak best
practices, rather than
open-ended reasoning.

\subsection{Heavy orchestration frameworks rejected}

Frameworks such as LangChain or LlamaIndex are often the default choice
for wiring a language model into an application: they bundle prompt
templates, retrieval, memory, and agent loops behind a single API, and
adopting one would have saved a little boilerplate around the call to
Ollama. They were rejected in favor of a direct, thin HTTP client that
posts to Ollama's own generation endpoint and normalizes every transport
failure into a single error type. The motivation goes beyond a general
dislike of dependencies: a framework interposes a layer of abstraction
between the request and the response that is harder to test, harder to
reason about when something goes wrong during a live demonstration, and
harder to explain precisely to a jury asking what the AI layer actually
does with a prompt. Section~\ref{principle:secure-by-design} depends on
being able to state exactly what goes into the model and what comes out of
it, since that is what keeps the no-PII guarantee on the prompt verifiable;
a direct call keeps that statement true by construction, whereas a
framework would require trusting that it adds nothing unexpected on either
side.

\subsection{Scoring logic kept out of the API routes}

The most direct way to expose a score over the API would have been to let
each route that needs one build the engine's input structure locally and
call the scoring functions itself, especially once a second consumer
needed the same translation: the recommendations endpoint requires the
same freshly computed \texttt{Assessment\allowbreak Result} as the score
endpoint itself, only to build a prompt from it instead of to serve it
directly. That path was rejected in favor of a single
bridge module that is the only place in the codebase allowed to translate
a stored assessment into the engine's input type and to call the engine.
Every route that needs a score goes through that one seam. The reason is
Section~\ref{principle:determinism}: two independent translations from the
database to the engine's input would eventually drift apart, for instance
one route dropping an incomplete answer differently from another, and the
platform would then be able to produce two different scores for the same
assessment depending on which endpoint computed it. A single bridge,
cross-checked in an end-to-end test against the engine called directly on
the same inputs, makes that kind of silent divergence structurally
impossible rather than merely unlikely.

\begin{table}[htbp]
  \centering
  \begin{tabular}{@{}p{0.3\textwidth}p{0.32\textwidth}p{0.18\textwidth}@{}}
    \toprule
    Choice made & Alternative rejected & Principle(s) \\
    \midrule
    Local inference via Ollama & Cloud LLM API (OpenAI, Anthropic, Mistral) & P2 \\
    Direct HTTP client to Ollama & Orchestration framework (LangChain, LlamaIndex) & P3 \\
    Single \texttt{scoring\_\allowbreak bridge.\allowbreak py} translation seam & Scoring logic duplicated across API routes & P1 \\
    \bottomrule
  \end{tabular}
  \caption{Summary of key design choices and rejected alternatives}
  \label{tab:rejected-alternatives}
\end{table}